% Ce chapitre repose sur :
%   DOC: s20, s29, s30
%   OBS: (aucun)
%   HYP: (aucun)
%   CHI: admissions-2024, capacite-annoncee, capacite-fonctionnelle, colonnes-part-doc, colonnes-part-hyp, colonnes-part-obs
%   CHI: colonnes-source-doc, colonnes-source-hyp, colonnes-source-obs, colonnes-source-registre, consultations-2024, consultations-produites
%   CHI: defaut-ages-incoherents, defaut-dates-aberrantes, defaut-factures-sans-pec, defaut-familles, defaut-rdv-doublon-creneau, defaut-taux-doublons
%   CHI: dim-patient-fiches-plusieurs-versions, dms-publie, ecart-annees-sources, examens-2024, facteur-cadrage-capacite-arrondi, facteur-cadrage-sejours-arrondi
%   CHI: factures-avec-creance, grappes-multiples, identifiants-distincts, interventions-2024, journees-2024, mouvements-lits-distincts
%   CHI: parametres-doc, parametres-fichiers, parametres-hyp, parametres-obs, parametres-part-doc, parametres-part-hyp
%   CHI: parametres-part-obs, parametres-total, part-sejours-depuis-urgences, passages-urgences-produits, patients-produits, periode-debut
%   CHI: periode-fin, periode-jours, personnes-estimees, prelevements-2024, sejours-produits, source-creances-lignes
%   CHI: source-encaissements-lignes, source-factures-lignes, source-lignes-facture-lignes, source-lignes-total, source-mouvements-lignes, source-passages-lignes
%   CHI: source-passages-urgences-lignes, source-patients-lignes, source-prises-en-charge-lignes, source-relances-lignes, source-rendez-vous-lignes, source-tables
%   CHI: tom-publie, urgences-borne-basse, urgences-borne-haute, urgences-milieu-fourchette, urgences-scenario-retenu, verifications-recueil
%   SER: (aucune)
%   SECTIONS: 6
\chapter{Conception du jeu de données}\label{chap:jeu-de-donnees}

Le jeu de données remplace les fichiers que le service ne peut pas confier. Il doit reprendre la
structure des écrans observés, des volumes plausibles pour l'hôpital et des défauts connus à
l'avance. La question traitée est double : comment le jeu est construit, et d'où vient chacune de
ses valeurs.

\section{Le modèle conceptuel}

Le modèle distingue la personne, ses identifiants et les versions de sa fiche (voir
figure~\ref{fig:diagramme-classes}). Une personne reçoit un ou plusieurs identifiants patient, et
c'est là que naît le problème d'identité du chapitre~\ref{chap:qualite}. À chaque identifiant se
rattachent des épisodes de soin de trois familles : consultation externe, passage aux urgences et
séjour. Un épisode donne lieu à une prise en charge, à une facture et, si elle reste impayée, à une
créance puis à des relances.

\begin{figure}[htbp]
\centering
\begin{tikzpicture}[node distance=13mm and 16mm]

\node[classeUML] (per) {\textbf{Personne}\nodepart{two}\emph{aucune clé au\\système d'origine}};
\node[classeUML, below=of per] (ipp) {\textbf{Identifiant patient}\nodepart{two}\texttt{n\_ipp}};
\node[classeUML, below=of ipp] (ver) {\textbf{Version de fiche}\nodepart{two}\texttt{nom}\\\texttt{sexe}\\\texttt{date\_naissance}\\\texttt{date\_extraction}};

\node[classeUML, right=of ipp] (epi) {\textbf{Épisode de soin}\nodepart{two}\texttt{n\_passage}\\\texttt{type\_passage}\\\texttt{date\_heure\_entree}\\\texttt{date\_heure\_sortie}};
\node[classeUML, above=22mm of epi] (rdv) {\textbf{Rendez-vous}\nodepart{two}\texttt{n\_rdv}\\\texttt{date\_rendez\_vous}\\\texttt{etat}};
\node[classeUML, below=of epi] (mvt) {\textbf{Mouvement}\nodepart{two}\texttt{n\_sejour}\\\texttt{n\_mutation}\\\texttt{mode\_admission}\\\texttt{mode\_sortie}};

\node[classeUML, right=of epi] (fac) {\textbf{Facture}\nodepart{two}\texttt{n\_facture}\\\texttt{montant\_total}\\\texttt{part\_organisme}\\\texttt{part\_patient}};
\node[classeUML, below=of fac] (cre) {\textbf{Créance}\nodepart{two}\texttt{n\_creance}\\\texttt{montant\_du}\\\texttt{montant\_restant}};
\node[classeUML, below=of cre] (rel) {\textbf{Relance}\nodepart{two}\texttt{n\_relance}\\\texttt{date\_relance}\\\texttt{canal}\\\texttt{resultat}};

% --- les associations, avec leurs cardinalités ---------------------------------
% Verticales : la cardinalité se pose à droite du trait, près de la classe qu'elle qualifie.
\draw[assocUML] (per) -- node[cardUML,right,pos=0.28] {1} node[cardUML,right,pos=0.78] {1..*} (ipp);
\draw[assocUML] (ipp) -- node[cardUML,right,pos=0.28] {1} node[cardUML,right,pos=0.78] {1..*} (ver);
\draw[assocUML] (epi) -- node[cardUML,right,pos=0.15] {1} node[cardUML,right,pos=0.85] {0..*} (mvt);
\draw[assocUML] (fac) -- node[cardUML,right,pos=0.15] {1} node[cardUML,right,pos=0.85] {0..1} (cre);
\draw[assocUML] (cre) -- node[cardUML,right,pos=0.15] {1} node[cardUML,right,pos=0.85] {0..*} (rel);
\draw[assocUML] (rdv) -- node[cardUML,right,pos=0.18] {0..1} node[cardUML,right,pos=0.82] {0..1} (epi);

% Horizontales : la cardinalité se pose au-dessus du trait.
\draw[assocUML] (ipp) -- node[cardUML,above,pos=0.18] {1} node[cardUML,above,pos=0.82] {0..*} (epi);
\draw[assocUML] (epi) -- node[cardUML,above,pos=0.18] {1} node[cardUML,above,pos=0.82] {0..*} (fac);

% En diagonale : le rendez-vous se rattache à l'identifiant, non à l'épisode. Le trait part du
% coin de la classe et monte vers la droite, si bien qu'il ne traverse aucune autre classe.
\draw[assocUML] (ipp.north east) -- node[cardUML,below right,pos=0.14] {1}
  node[cardUML,above left,pos=0.94] {0..*} (rdv.west);
\end{tikzpicture}
\caption{Les entités du domaine et leurs cardinalités}\label{fig:diagramme-classes}
\end{figure}

Les trois niveaux de l'identité donnent trois décomptes différents. La couche produite porte
\chiffre{patients-produits} versions de fiche, réparties sur \chiffre{identifiants-distincts}
identifiants, dont \chiffre{dim-patient-fiches-plusieurs-versions} portent plus d'une version. Ces
identifiants désignent \chiffre{personnes-estimees} personnes estimées, car le jeu injecte des
doublons d'identité : \chiffre{grappes-multiples} grappes réunissent plus d'un identifiant. Le
registre des chiffres porte une entrée par décompte, et la note de chacune renvoie aux deux autres.

Deux autres choix du modèle servent l'analyse. Le rendez-vous est une entité distincte de
l'épisode, et son lien à l'épisode est facultatif. Un rendez-vous non honoré n'a donc pas
d'épisode, et c'est ce qui rend l'absentéisme mesurable. La facture est elle aussi facultative,
parce qu'un épisode non facturé est un défaut réel du service que la Cour des comptes relève
\cite{s20}. Le chapitre~\ref{chap:activite} en fait un indicateur.

\section{La couche source}

La couche source reproduit ce qu'une extraction du logiciel rendrait : \chiffre{source-tables}
tables, une par flux, et \chiffre{source-lignes-total} lignes sur la période (voir
tableau~\ref{tab:tables-source}).

\begin{table}[htbp]
\centering
\caption{Les tables de la couche source, leur grain et leur volume sur la période}\label{tab:tables-source}
\begin{tabular}{@{}p{3.7cm}p{5.9cm}r@{}}
\toprule
Table & Grain — une ligne par & Lignes \\
\midrule
\texttt{patients} & fiche patient, dans son état à l'extraction & \chiffre{source-patients-lignes} \\
\texttt{rendez\_vous} & rendez-vous donné & \chiffre{source-rendez-vous-lignes} \\
\texttt{passages} & épisode de soin, toutes familles confondues & \chiffre{source-passages-lignes} \\
\texttt{passages\_urgences} & passage aux urgences, avec son tri & \chiffre{source-passages-urgences-lignes} \\
\texttt{mouvements} & mouvement de séjour : admission, mutation ou sortie & \chiffre{source-mouvements-lignes} \\
\texttt{prises\_en\_charge} & prise en charge d'un épisode par un organisme & \chiffre{source-prises-en-charge-lignes} \\
\texttt{factures} & facture émise & \chiffre{source-factures-lignes} \\
\texttt{lignes\_facture} & ligne d'acte d'une facture & \chiffre{source-lignes-facture-lignes} \\
\texttt{encaissements} & encaissement rattaché à une facture & \chiffre{source-encaissements-lignes} \\
\texttt{creances} & créance née d'une facture non soldée & \chiffre{source-creances-lignes} \\
\texttt{relances} & relance émise sur une créance & \chiffre{source-relances-lignes} \\
\midrule
\textbf{Total} & & \textbf{\chiffre{source-lignes-total}} \\
\bottomrule
\end{tabular}
\end{table}

Les clés suivent le système d'origine. L'identifiant patient relie les fiches aux épisodes. Le
numéro de facture relie la facture à ses lignes, à ses encaissements et à sa créance. Le numéro de
séjour relie les mouvements entre eux. Aucune clé de substitution n'est ajoutée à ce niveau, et les
défauts de l'extraction restent en place.

Deux traits de cette couche pèsent sur les calculs du chapitre~\ref{chap:activite}. Seules
\chiffre{factures-avec-creance} des \chiffre{source-factures-lignes} factures donnent lieu à une
créance : un taux de recouvrement change donc de plusieurs dizaines de points selon la population
retenue. La table des passages porte tous les épisodes, distingués par un type, et celle des
passages aux urgences ajoute une ligne par passage d'urgence, avec son niveau de tri. Ses
\chiffre{source-passages-urgences-lignes} lignes sont exactement les passages de type urgence de la
première, et additionner les deux tables compterait deux fois la même activité.

\section{La provenance des colonnes et des paramètres}

Chaque colonne de la couche source et chaque paramètre du générateur portent une étiquette qui dit
d'où ils viennent (voir tableau~\ref{tab:structure-et-valeurs}). Un registre unique décrit les
\chiffre{colonnes-source-registre} colonnes. Il dit pour chacune si le champ a été vu à l'écran,
déduit d'une source documentaire ou posé sans preuve externe. La définition de schéma et le
dictionnaire des données en sont produits mécaniquement (annexe~\ref{annexe:dictionnaire}), et un
contrôle bloquant vérifie que la base porte exactement ces colonnes. La correspondance avec les
écrans, colonne par colonne, et le décompte au catalogue de la base sont donnés en
section~\ref{sec:complements-provenance}.

\begin{table}[htbp]
\centering
\caption{Provenance de la structure et provenance des valeurs, côte à côte}\label{tab:structure-et-valeurs}
\begin{tabular}{@{}p{2.2cm}rrrr@{}}
\toprule
Étiquette & \multicolumn{2}{c}{Colonnes (structure)} & \multicolumn{2}{c}{Paramètres (valeurs)} \\
\midrule
\texttt{OBS} & \chiffre{colonnes-source-obs} & \chiffre{colonnes-part-obs}~\% & \chiffre{parametres-obs} & \chiffre{parametres-part-obs}~\% \\
\texttt{DOC} & \chiffre{colonnes-source-doc} & \chiffre{colonnes-part-doc}~\% & \chiffre{parametres-doc} & \chiffre{parametres-part-doc}~\% \\
\texttt{HYP} & \chiffre{colonnes-source-hyp} & \chiffre{colonnes-part-hyp}~\% & \chiffre{parametres-hyp} & \chiffre{parametres-part-hyp}~\% \\
\midrule
\textbf{Total} & \textbf{\chiffre{colonnes-source-registre}} & & \textbf{\chiffre{parametres-total}} & \\
\bottomrule
\end{tabular}
\end{table}

Le tableau montre que les valeurs sont bien plus hypothétiques que la structure qui les porte. Les
colonnes vues à l'écran font \chiffre{colonnes-part-obs}~\% des colonnes, contre
\chiffre{parametres-part-obs}~\% pour les paramètres observés. À l'inverse,
\chiffre{parametres-part-hyp}~\% des paramètres sont posés sans preuve externe, contre
\chiffre{colonnes-part-hyp}~\% des colonnes. Les paramètres sont comptés sur les
\chiffre{parametres-fichiers} fichiers de configuration du générateur, et ce sont eux qui produisent
les nombres du chapitre~\ref{chap:activite}.

\section{La volumétrie}

La volumétrie du jeu est relevée pour l'hôpital lui-même dans le recueil statistique national, puis
mise à l'échelle de la période. La période générée court du \chiffre{periode-debut} au
\chiffre{periode-fin}, soit \chiffre{periode-jours} jours. Pour l'exercice 2024, le recueil publie
une capacité fonctionnelle de \chiffre{capacite-fonctionnelle} lits, \chiffre{journees-2024}
journées d'hospitalisation et \chiffre{admissions-2024} admissions \cite{s30}. Il publie aussi
\chiffre{consultations-2024} consultations spécialisées externes, \chiffre{prelevements-2024}
prélèvements et \chiffre{examens-2024} examens de laboratoire. La couche produite porte en retour
\chiffre{sejours-produits} séjours, \chiffre{consultations-produites} consultations et
\chiffre{passages-urgences-produits} passages aux urgences.

Le cadrage initial avait dérivé la volumétrie d'une capacité annoncée et d'indicateurs nationaux,
et se trompait d'un facteur \chiffre{facteur-cadrage-capacite-arrondi} sur la capacité et \chiffre{facteur-cadrage-sejours-arrondi} sur les séjours \cite{s30}. Il confondait la capacité annoncée par la presse,
\chiffre{capacite-annoncee} lits \cite{s29}, avec la capacité fonctionnelle, seule publiée par
établissement. Il appliquait aussi des valeurs nationales à un hôpital qui publie les siennes : un
taux d'occupation de \chiffre{tom-publie}~\% et une durée moyenne de séjour de \chiffre{dms-publie}
jours. Compter les lits distincts des mouvements rendrait \chiffre{mouvements-lits-distincts} lits,
la valeur publiée. Ce décompte n'est pas retenu, car un lit resté vide n'y apparaît pas : il donne
un plancher, et non une capacité.

Le recueil porte \chiffre{interventions-2024} intervention chirurgicale pour l'hôpital \cite{s30},
et ce zéro vient d'une absence de ligne plutôt que d'un zéro imprimé. La Cour des comptes décrit de
même, à \chiffre{ecart-annees-sources} ans d'écart, une gynéco-obstétrique qui fonctionne comme maison d'accouchement, sans
activité chirurgicale, malgré un bloc opératoire équipé \cite{s20}. Les \chiffre{verifications-recueil} vérifications qui fondent cette lecture du
recueil sont données en section~\ref{sec:complements-chirurgie}.

\section{La sensibilité aux passages aux urgences}\label{sec:sensibilite}

Le volume des passages aux urgences est la seule valeur absolue du modèle qu'aucune source ne
publie, et son choix ne touche aucun taux ni aucun délai. Le recueil statistique ne donne ce volume
à aucun niveau géographique. La fourchette retenue est donc construite, et ses deux bornes n'ont pas
la même solidité.

La borne haute, \chiffre{urgences-borne-haute} passages par jour, transpose l'activité d'un hôpital
voisin rapportée à sa capacité litière \cite{s20}. Le même rapport la contredit : les urgences de
l'hôpital étudié n'assuraient pas la garde, faute de généralistes affectés. Le chiffre de référence
précède en outre la généralisation de l'assurance maladie obligatoire, dont l'effet sur le recours
aux urgences publiques n'est pas connu. La borne basse, \chiffre{urgences-borne-basse} passages par
jour, combine la part des séjours venus des urgences et un taux d'hospitalisation aux urgences, deux
ratios qui ne sont pas sourcés.

Le scénario retenu, \chiffre{urgences-scenario-retenu} passages par jour, est placé sous le milieu
de la fourchette, qui vaut \chiffre{urgences-milieu-fourchette}. La borne haute repose en effet sur
une transposition que la seule description du service contredit. La borne basse souffre seulement de ratios non sourcés.

Le scénario fixe le volume des passages aux urgences, \chiffre{passages-urgences-produits} sur la
période, et ce qui en dérive. Il fixe ainsi la taille de la table des passages aux urgences. Il
fixe aussi le taux d'hospitalisation aux urgences, qui rapporte la part des séjours venus des urgences,
\chiffre{part-sejours-depuis-urgences}, aux passages du scénario.

Les séjours restent ancrés sur les
admissions relevées, les délais sont tirés de distributions et le rapprochement dépend du nombre de
fiches patient. Un taux voit son numérateur et son dénominateur changer ensemble. Aucun résultat
exprimé en taux ou en délai ne dépend donc du scénario, et les volumes de passages aux urgences sont
toujours présentés comme une hypothèse.

\section{Les distributions et les défauts injectés}

Les valeurs du générateur suivent des distributions sourcées ou posées, et des défauts y sont
injectés pour éprouver la chaîne. Chaque distribution porte une étiquette de provenance, et la
liste complète est donnée en section~\ref{sec:complements-distributions}. Les distributions de
l'identité et de la couverture sont posées, celles de la gravité aux urgences sont documentées.

Une étiquette dit d'où vient une valeur, et non qu'elle s'applique au service. La répartition par
niveau de tri reproduit ainsi une échelle française publiée, sans montrer qu'un service marocain
trie de la même façon. Chaque paramètre porte aussi, dans sa configuration, la phrase qui dit ce
qui serait faux si sa valeur l'était.

Les défauts sont injectés parce qu'un jeu propre ne permet de tester aucun traitement de qualité.
Leur liste est enregistrée, si bien que ce qu'un traitement devrait trouver est connu.
Le jeu porte \chiffre{defaut-familles} familles de défauts. Les principales portent un taux
explicite : \chiffre{defaut-taux-doublons} pour les doublons d'identité,
\chiffre{defaut-rdv-doublon-creneau} pour les rendez-vous en doublon de créneau,
\chiffre{defaut-factures-sans-pec} pour les factures sans prise en charge,
\chiffre{defaut-dates-aberrantes} pour les dates aberrantes et \chiffre{defaut-ages-incoherents}
pour les âges incohérents. S'y ajoutent des champs manquants, à un taux propre à chaque colonne, et
des défauts de casse, d'accents et d'espaces.

Ces défauts forment trois classes, selon ce qui les révèle. Une date au format invalide se voit sur
une ligne seule, et la ligne part en quarantaine dès le chargement (section~\ref{sec:quarantaine}).
Un doublon d'identité ne se voit qu'en comparant deux fiches, valides l'une et l'autre. D'autres
défauts ne se voient qu'en joignant deux tables.

Une quatrième forme échappe aux trois : un champ renseigné et faux. Au poste, le numéro de
téléphone est parfois rempli d'un numéro inventé lorsque la personne ne connaît pas le sien
(chapitre~\ref{chap:systeme}). Le générateur reproduit ce défaut : une part des numéros produits a
une forme valide et ne désigne personne. Aucun contrôle de forme ne peut l'attraper.

\section*{Synthèse}

Le jeu de données reprend la structure des écrans observés et les volumes publiés pour l'hôpital.
Ses colonnes sont souvent observées, ses paramètres le plus souvent posés. Le volume des urgences
est la seule valeur absolue sans source, et il ne touche aucun taux. Les défauts injectés sont
connus, ce qui permet de mesurer ce que la chaîne en retrouve. Le chapitre~\ref{chap:architecture}
décrit la chaîne qui charge et transforme ce jeu.
